\section{Introduction }
In 1965 Gordon Moore postulated Moores law, that the number of transistors on an integrated circuit doubles every 2 years \cite{4785860}. 
This law held for 50 where the transistor count increased form 2308 to 58,2 billion \cite{owid-moores-law}. 
Recently the increase in transistors is slowing down and Moores law is predicted to break for among other reasons manufacturing and material limits and the heat wall \cite{Markov_2014}.
To overcome these fundamental limits magnetic skyrmions have emerged as a promising alternative to electron based electronics. 
Spintronics is a promising candidate for high performance and low energy consumption computing and data storage \cite{Zhang_2023}. 
Magnetic skyrmions are topologically protected spin textures, the formation of magnetic skyrmions is generally attributed to the competition between ferromagnetic exchange interactions and chiral spin-orbit coupling terms such as the Dzyaloshinskii Moriya interaction (DMI) \cite{Dzyaloshinsky1958ATT} \cite{PhysRev.120.91}. 
In this work, we investigate the emergence and stability of skyrmionic configurations in a two-dimensional classical spin lattice described by a Heisenberg Dzyaloshinskii Moriya Hamiltonian. 
The system consists of unit vector spins arranged on a square lattice with periodic boundary conditions. 
To explore the phase space of the model, Monte Carlo simulations are performed over a range of Dzyaloshinskii Moriya interaction strengths and magnetic field values. 
For each parameter set, the system is annealed through a sequence of inverse temperatures until a stable configuration is obtained. 
The resulting spin textures are analyzed using a discretized topological charge, computed from the solid angles subtended by neighboring spins on lattice plaquettes \cite{BERG1981412}. This discretized winding number provides a robust measure of the skyrmion number and allows the identification of topologically non trivial states on the lattice.
